% ------------------------------------------------------------
% Abstract
% ------------------------------------------------------------
\begin{abstract}
Deep neural networks denoise low-dose CT (LDCT) with impressive PSNR and SSIM, yet recent benchmarking indicates that six years of architectural progress on these metrics has been ``marginal at best''~\cite{eulig2024}. We argue that the more consequential failure axis is \emph{physical} --- and that it is not specific to any particular network. Denoisers trained with image-appearance objectives (pixel MSE, but equally SSIM-, perceptual-, and adversarially-driven losses) optimize how images \emph{look}, not how noise \emph{behaves}, and therefore generically distort the noise power spectrum (NPS) and the Hounsfield-unit (HU) calibration of their output. Both distortions carry potential clinical cost: HU errors can corrupt attenuation-based decisions and quantitative pipelines (lesion characterization, calcium scoring, radiomics, radiotherapy planning), while a distorted NPS changes the noise texture radiologists are trained on and can degrade low-contrast lesion detectability. Auditing a benchmark-faithful, classically strong RED-CNN as a representative testbed, we find that it leaves 16\% of the noise magnitude unremoved overall (26\% in the abdomen), reproduces the abdominal noise spectral shape with a correlation of only 0.51, and introduces tissue-dependent HU biases up to $-53$\hu{} in chest bone --- all invisible to PSNR and SSIM. We propose a physics-faithful framework that combines two lightweight physics losses (radial log-NPS matching and per-tissue HU-bin bias) with an architectural contribution: a \textbf{DC-preserving, anatomy-adaptive radial spectral head} --- a ${\approx}5$k-parameter module with a trunk-agnostic interface that reshapes the removed-noise residual in the Fourier domain under provable per-image guarantees (exact DC preservation of the residual correction, frozen near-DC gains, isotropic radial parameterization, bounded adaptive offsets). At strictly matched training budgets and loss weights, adding only the head improves every measured physics metric --- NPS error $-12\%$ overall and $-19\%$ in the chest (mean over three seeds), noise-magnitude ratio 0.93 vs.\ 0.91, chest soft-tissue and bone biases substantially reduced --- at equal classical image quality, winning the pre-specified primary spectral metric in 30/30 patient--seed comparisons (10 patients $\times$ 3 seeds; exact two-sided Wilcoxon $p = 0.002$ within each seed). We further show that physics losses alone \emph{saturate}: doubling their weight without the head degrades both physics and image quality, identifying the head as the mechanism that converts stronger physics incentives into spectral fidelity rather than instability. Multi-seed replication exposes an additional failure of appearance-only training: its physics is seed-unstable ($\pm 13\%$ relative NPS spread at pinned PSNR/SSIM), and the physics constraints cut this variance by $2.5$--$4\times$. A cross-trunk replication on a ResNet trunk shows the physics losses transfer unchanged (NPS error $-39\%$ with no re-tuning), while the head's quantitative benefit is trunk-dependent --- the interface is trunk-agnostic, the learned equilibrium is not --- a boundary we measure, diagnose via the head's interpretable gain curves, and report explicitly. The adaptive component learns quantifiably anatomy-specific corrections (chest/abdomen gain-curve separation ${\approx}4\times$ the within-anatomy spread) without sacrificing any guarantee, and yields directly interpretable, publishable gain curves.
\end{abstract}
\vspace{0.5em}
